\chapter{ARS MORIENDI}\label{ars-moriendi}

\hfill\break

I marziani, mi spiegò Jason, non sono quel popolo semplice, pacifico e
bucolico come Wun ci aveva fatto (o lasciato) credere.

Era vero che non erano particolarmente bellicosi - le Cinque Repubbliche
avevano risolto le loro controversie quasi un millennio prima - ed erano
``bucolici'' nel senso che avevano stanziato la maggior parte delle loro
risorse nell'agricoltura. Ma non erano ``semplici'' in alcun senso
possibile della parola. Erano, come sottolineò Jase, maestri nell'arte
della biologia sintetica. La loro civiltà era basata su quello. Avevamo
costruito loro un pianeta abitabile con strumenti biotecnologici, e non
c'è mai stata una generazione marziana che non abbia capito la funzione
e i possibili utilizzi del DNA.

Se la loro tecnologia su vasta scala poteva sembrare grezza - la
navicella di Wun, per esempio, era quasi primitiva, un cannone
newtoniano - lo si doveva alle risorse naturali profondamente limitate.
Marte era un mondo senza petrolio né carbone a sostenere un ecosistema
fragile e povero di acqua e azoto. Un tessuto industriale dissoluto e
lussureggiante come il nostro non sarebbe mai potuto esistere sul
pianeta di Wun. Su Marte, gran parte degli sforzi dei primi umani erano
stati volti a produrre una quantità di cibo sufficiente a una
popolazione rigorosamente controllata. La biotecnologia era mirabilmente
utile allo scopo. Le industrie con le loro ciminiere no.

«Te l'ha detto Wun?» chiesi, mentre pioveva scrosciava senza
interruzione e il pomeriggio scivolava via.

«Si fidava di me, sì, anche se la maggior parte delle cose che mi ha
detto erano sottintese negli archivi.»

La luce color ruggine proveniente dalla finestra si rifletteva negli
occhi ciechi e modificati di Jason.

«Ma potrebbe aver mentito.»

«Non so se abbia mai \emph{mentito}, Tyler. Era solo un po' tirchio di
realtà.»

I microscopici replicatori che Wun aveva portato sulla Terra erano
biologia sintetica all'avanguardia. Erano in grado di fare tutto ciò che
Wun aveva assicurato che avrebbero fatto. In effetti, erano molto più
complicati di quanto abbia mai voluto ammettere Wun.

Tra le varie funzioni nascoste dei replicatori c'era un sottocanale per
la comunicazione tra di loro e con il punto di origine. Wun non aveva
spiegato se si trattasse di un normale segnale radio a banda stretta o
di una trasmissione più insolita dal punto di vista tecnologico - Jase
propendeva per la seconda ipotesi. In ogni caso, richiedeva un
ricevitore molto più avanzato di quello che avremmo potuto mai costruire
sulla Terra. Aveva bisogno, come aveva detto Wun, di un ricevitore
biologico. Un sistema nervoso umano modificato.

ooo

«Ti sei offerto \emph{volontario}?»

«Lo avrei fatto. Se qualcuno me lo avesse chiesto. Ma l'unico motivo per
cui Wun si fidava di me era che temeva per la propria vita dal primo
giorno in cui era arrivato sulla Terra. Non aveva dubbi sulla venalità
umana o sulle politiche di potere. Gli serviva qualcuno di cui si fidava
che fosse disponibile a prendere in custodia la sua farmacopea nel caso
gli fosse capitato qualcosa. Qualcuno che ne capisse lo scopo. Non mi ha
mai proposto di diventare un ricevitore. La modificazione funziona
soltanto su un Quarto - ricordi cosa ti ho detto? Il trattamento di
longevità è una piattaforma. Fa funzionare altre applicazioni. Questa è
una di quelle.»

«Ti sei causato da solo tutto questo \emph{di proposito}?»

«Mi sono iniettato la sostanza dopo la sua morte. Non è stato traumatico
e non ho subito alcun effetto collaterale immediato. Ricorda, Tyler, non
c'era modo che una comunicazione dei replicatori potesse attraversare la
membrana Spin senza venire deteriorata. Ciò che mi sono dato era
un'abilità \emph{latente}.»

«Allora, perché farlo?»

«Perché non volevo morire nell'ignoranza. Tutti quanti abbiamo dato per
scontato che, una volta finito lo Spin, saremmo morti nel giro di pochi
giorni o di ore. L'unico vantaggio garantito dalla modificazione di Wun
era che in quegli ultimi giorni oppure ore, finché fossi sopravvissuto,
sarei stato in stretto contatto con una banca dati grande quasi quanto
la galassia stessa. Avrei saputo più di quanto chiunque sulla Terra
avrebbe \emph{potuto} sapere, chi erano gli Ipotetici e perché ci hanno
fatto tutto questo.»

Pensai: ``E lo sai, adesso?'' Ma forse lo sapeva davvero. Forse era
proprio questo che voleva raccontare prima di perdere la facoltà della
parola, la ragione per cui voleva che facessi una registrazione. «Wun
sapeva che avresti potuto farlo?»

«No, e dubito che avrebbe approvato\ldots{} sebbene anche lui stesse
facendo funzionare la stessa applicazione.»

«Davvero? Non lo dava a vedere.»

«Non ce n'era ragione. Ricorda: ciò che sta capitando a me - al mio
corpo, al mio cervello - non è l'applicazione.» Rivolse gli occhi senza
vista verso di me. «Questo è un guasto.»

ooo

I replicatori erano stati lanciati dalla Terra e avevano prosperato nel
sistema solare esterno, lontano dal Sole. (E se gli Ipotetici se ne
fossero accorti e avessero dato la colpa alla Terra di ciò che in realtà
era un intervento marziano? Era forse proprio questo, come aveva
lasciato intendere E.D., ciò che i subdoli marziani avevano in mente fin
dall'inizio? Jason non lo disse - presunsi che non lo sapeva.)

Con il tempo i replicatori si erano diffusi sulle stelle più vicine e
oltre\ldots{} alla fine, molto oltre. Le colonie dei replicatori erano
invisibili a distanze astronomiche, ma se li avessimo potuti mappare su
una griglia del nostro locale circondario stellare, avremmo potuto
vedere una nuvola di replicatori in continua espansione, un'esplosione
glacialmente lenta di vita artificiale.

I replicatori non erano immortali. Come entità individuali vivevano, si
riproducevano e alla fine morivano. Ciò che restava era la rete che
avevano costruito: una barriera corallina di nodi recintati e
interconnessi in cui i nuovi dati si accumulavano e venivano drenati
verso il punto d'origine della rete.

«L'ultima volta che abbiamo parlato,» ricordai a Jase, «mi avevi detto
che c'era un problema, che la popolazione dei replicatori stava
regredendo.»

«Avevano incontrato qualcosa che nessuno aveva previsto.»

«Cos'era, Jase?»

Rimase in silenzio per alcuni attimi, come per raccogliere le idee.

«Supponevamo,» continuò, «che con il lancio dei replicatori stessimo
introducendo nell'universo qualcosa di nuovo, un tipo del tutto inedito
di vita artificiale. Quella supposizione era ingenua. Noi - esseri
umani, terrestri o marziani - non siamo stati i primi esseri senzienti
evoluti nella nostra galassia. Anzi! Infatti, non abbiamo nulla di
particolarmente inusuale. In teoria, tutto ciò che abbiamo fatto nella
nostra breve storia è già stato fatto prima, da qualcun altro.»

«Mi stai dicendo che i replicatori si sono imbattuti in altri
replicatori?»

«In \emph{un'ecologia} di replicatori. Le stelle sono una giungla,
Tyler. Contengono più vita di quanto abbiamo mai immaginato.»

Provai a visualizzare il processo mentre Jason me lo descriveva.

Molto distante dalla Terra isolata dallo Spin, ben oltre il sistema
solare - così in profondità nello spazio che il Sole stesso è solo una
stella qualsiasi in un cielo affollato - un seme di replicatore si era
posato su un polveroso frammento di ghiaccio e aveva iniziato a
riprodursi. Aveva iniziato lo stesso ciclo di crescita,
specializzazione, osservazione, comunicazione e riproduzione che aveva
avuto luogo innumerevoli volte durante le lente migrazioni dei suoi
antenati. Forse avrebbe potuto raggiungere la maturità; forse avrebbe
potuto addirittura cominciare a emettere microvortici di dati; ma quella
volta il ciclo si era interrotto.

Qualcosa aveva percepito la presenza dei replicatori. Qualcosa che aveva
fame.

Il predatore (mi spiegò Jase) era un altro tipo di sistema semiorganico
autocatalitico di retroazione - un'altra colonia di meccanismi cellulari
autoriproduttivi, tanto artificiale quanto biologica - e il predatore
era collegato alla sua stessa rete, questa in particolare più vecchia e
vasta di qualsiasi altra cosa avessero avuto il tempo di costruire i
replicatori terrestri durante il loro esodo dalla Terra. Il predatore
era di gran lunga più evoluto della sua preda: le sue subroutine per la
ricerca di nutritivi e l'utilizzo delle risorse erano state affinate nel
corso di miliardi di anni. La colonia replicatrice terrestre, cieca e
incapace di fuggire, era stata prontamente divorata.

Ma ``divorata'' in quel caso assumeva un significato speciale. Il
predatore voleva di più delle sofisticate molecole di carbonio di cui la
forma matura del replicatore era composta, per quanto utili potessero
essere. Ben più interessante, per il predatore, era il
\emph{significato} del replicatore, le funzioni e le strategie scritte
nei suoi modelli riproduttivi. Adottava da questi ultimi ciò che
considerava di potenziale valore; poi riorganizzava e sfruttava la
colonia dei replicatori per i propri scopi. La colonia non moriva ma
veniva assorbita, ontologicamente divorata, inglobata con i suoi
confratelli all'interno di una gerarchia interstellare più vasta, più
complessa e molto, molto più vecchia.

Non sarebbe stata né la prima né l'ultima volta che un dispositivo del
genere veniva assorbito.

«Le reti dei replicatori,» disse Jason, «sono una delle cose che le
civiltà senzienti tendono a produrre. Data l'inerente difficoltà del
viaggio a velocità subluce come mezzo d'esplorazione della galassia, la
maggior parte delle culture tecnologiche alla fine si accontenta di
usare un'estesa rete di macchine di von Neumann, in sostanza la
definizione dei replicatori, priva di costi di mantenimento e in grado
di generare un effetto a cascata d'informazioni scientifiche, che si
espande esponenzialmente nel corso dei tempi storici.»

«Va bene,» gli risposi, «ho capito. I replicatori marziani non sono
unici. Si sono imbattuti in quella che tu chiami un'ecologia\ldots»

«Sì, un'ecologia di macchine von Neumann.» (In nome del matematico del
ventesimo secolo John von Neumann, che per primo suggerì la possibilità
di macchine autoreplicanti.)

«Un'ecologia di von Neumann, e ne sono stati assorbiti. Ma questo non ci
dice niente sugli Ipotetici o sullo Spin.»

Jason contrasse le labbra con impazienza. «Tyler, no. Non capisci. Gli
Ipotetici sono l'ecologia di von Neumann. Sono la stessa cosa.»

ooo

A questo punto dovetti fare un passo indietro e riconsiderare con
precisione chi era quello nella stanza con me.

Sembrava Jase. Ma tutto ciò che aveva detto lo metteva in dubbio.

«Stai per caso comunicando con questa\ldots{} entità? Voglio dire ora?
Mentre parliamo?»

«Non so se si possa definire comunicazione. La comunicazione funziona a
due sensi. Questa no, non nel senso che intendi tu. E una vera e propria
comunicazione non sarebbe così sconvolgente. Questa lo è. Soprattutto di
notte. L'impulso di dati in arrivo è moderato durante le ore diurne,
presumibilmente perché le radiazioni solari fanno perdere il segnale.»

«Di notte il segnale è più forte?»

«Forse anche la parola ``segnale'' è fuorviante. Un segnale è ciò che i
replicatori originali erano progettati a trasmettere. Quello che ricevo
io arriva sulla stessa onda portante, e convoglia informazioni, ma è
attivo, non passivo. Cerca di fare a me quello che ha già fatto a ogni
altro nodo della rete. In pratica, Ty, sta cercando di acquisire e
riprogrammare il mio sistema nervoso.»

Quindi c'era una terza entità nella stanza. Io, Jase\ldots{} e gli
Ipotetici, che lo stavano divorando vivo.

«Possono farlo? Riprogrammare il tuo sistema nervoso intendo?»

«Non in maniera efficace, no. Per loro sono solo l'ennesimo nodo nella
rete dei replicatori. La biotecnologia che mi sono iniettato è sensibile
alla loro manipolazione, ma non nei modi che si aspettano. Siccome non
mi percepiscono come entità biologica, l'unica cosa che possono fare è
uccidermi.»

«Esiste un modo per schermare questo segnale, per creare
un'interferenza?»

«Non che io sappia. Se i marziani erano in possesso di tale tecnica,
hanno omesso di includere l'informazione nei loro archivi.»

La finestra della camera di Jason guardava a ovest. Il bagliore rosato
che stava entrando ora nella stanza era il Sole che svaniva, oscurato
dalle nuvole.

«Ma adesso sono insieme a te. Ti parlano.»

«Loro. Essa. Ci vorrebbe un pronome migliore. L'intera ecologia di von
Neumann è un'entità singola. Sviluppa da sola i propri pensieri lenti ed
elabora i propri piani. Ma molte delle sue innumerevoli parti sono anche
individui autonomi, spesso in competizione tra loro, agiscono più
rapidamente rispetto alla rete nella sua totalità e sono di gran lunga
più intelligenti di ogni singolo essere umano. La membrana dello Spin,
per esempio\ldots»

«La membrana dello Spin è un \emph{individuo}?»

«Per molti aspetti, sì. I suoi obiettivi finali derivano dalla rete, ma
valuta gli eventi e fa scelte autonome. È più complesso di quanto ci
siamo mai sognati, Ty. Abbiamo dato per scontato che la membrana fosse
\emph{accesa o spenta}, come l'interruttore della luce, come un codice
binario. Non è così. Ha \emph{molte} posizioni. \emph{Molti} scopi.
Molti gradi di permeabilità, per esempio. Da anni sappiamo che può far
transitare una navicella spaziale e respingere un asteroide. Ma possiede
anche capacità più ingegnose. Ecco perché negli ultimi giorni non siamo
stati sopraffatti dalle radiazioni solari. La membrana ci sta ancora
offrendo un certo livello di protezione.»

«Non conosco i numeri delle vittime, Jase, ma solo in questa città ci
devono essere migliaia di persone che hanno perso la famiglia da quando
lo Spin si è fermato. Sarei molto restio a dire loro che sono stati
``protetti''.»

«Ma lo sono. In generale sì, anche se non singolarmente. La membrana
dello Spin non è Dio - non può vedere la caduta di un passero. Può,
tuttavia, impedire che il passero venga cotto da una luce ultravioletta
letale.»

«A che scopo?»

A queste parole si imbronciò. «Non riesco a capirlo di preciso,»
cominciò, «o forse non riesco a tradurlo\ldots»

Bussarono alla porta. Carol entrò con un mucchio di biancheria. Spensi
il registratore e lo misi via. L'espressione di Carol era lugubre.

«Lenzuola pulite?» chiesi.

«Fasce» rispose brusca. La biancheria era stata tagliata in strisce.
«Per quando iniziano le convulsioni.»

Annuì in direzione della finestra, la luce del giorno si allungava.

«Grazie» disse Jason dolcemente. «Tyler, se hai bisogno di una pausa
questo è il momento giusto. Ma non stare via troppo a lungo.»

ooo

Andai a controllare Diane che era in una fase d'intervallo e dormiva.
Pensai al farmaco marziano che le avevo somministrato (il ``Quattro di
base'', come l'aveva chiamato Jase), molecole semintelligenti sul punto
di dare battaglia al carico enorme di batteri della sindrome da
deperimento cardiovascolare, microscopici battaglioni radunati per
riparare e ricostruire il corpo di Diane, a patto che non fosse troppo
debole per sopportare lo stress della trasformazione.

Le baciai la fronte e le sussurrai parole dolci che probabilmente non
riuscì a sentire. Poi lasciai la camera e scesi giù dalle scale, e
quindi uscii, sul prato della Grande Casa, per ritagliare un istante
solo per me.

La pioggia era finalmente cessata - all'improvviso, del tutto - e l'aria
era più fresca di quanto lo fosse stata per l'intero giorno. Allo zenit
il cielo era blu scuro. Alcuni brandelli di nuvole di tempesta
ammantavano il Sole mostruoso nel punto in cui toccava l'orizzonte
occidentale. Gocce di pioggia tremolavano su ogni filo d'erba, minuscole
perle ambrate.

Jason aveva ammesso che stava morendo. Fu solo allora che cominciai ad
ammetterlo a me stesso.

In quanto medico, avevo già visto molta più morte della maggior parte
delle altre persone. Sapevo come moriva la gente. Sapevo che la famosa
storia su come affrontiamo la morte - negazione, rabbia, accettazione -
era al massimo una rozza generalizzazione. Certe emozioni possono
svilupparsi in pochi secondi o non farlo mai; la morte può sconfiggerle
in qualsiasi momento. Molta gente in realtà non si è trovata a dover
affrontare la morte; il loro trapasso è arrivato senza preavviso, una
rottura dell'aorta o una decisione sbagliata a un incrocio trafficato.

Ma Jase sapeva di stare morendo. E io ero sbalordito dalla calma
ultraterrena con cui sembrava averlo accettato, finché mi resi conto che
anche la sua morte era un'ambizione realizzata. Era sul punto di capire
ciò che aveva tentato di comprendere per tutta la vita: il significato
dello Spin e il ruolo che vi aveva giocato l'umanità - il \emph{suo}
ruolo, dal momento che era stato essenziale nel lancio dei replicatori.

Era come se avesse teso le mani fino a toccare le stelle.

E loro lo avevano \emph{toccato} in risposta. Le stelle lo stavano
ammazzando. Ma stava morendo in uno stato di grazia.

ooo

«Dobbiamo sbrigarci. Ormai è quasi buio, no?» mi chiese Jason.

Carol se n'era andata ad accendere delle candele per tutta la casa.

«Quasi» gli risposi.

«E la pioggia è finita. O almeno, non riesco a sentirla.»

«Anche la temperatura sta scendendo. Vuoi che apra la finestra?»

«Sì, per favore. E il registratore, lo hai riacceso?»

«Sì, è in funzione.» Alzai di qualche centimetro la vecchia cornice
della finestra e l'aria fresca entrò nella stanza.

«Stavamo parlando degli Ipotetici\ldots»

«Sì.» Silenzio. «Jase? Sei ancora con me?»

«Sento il vento. Sento la tua voce. Sento\ldots»

«Jason?»

«Scusami\ldots{} non farci caso, Ty. Ora come ora mi distraggo
facilmente. Io\ldots{} \emph{oh}!»

Le braccia e le gambe strattonarono le fasce con cui Carol lo aveva
bloccato al letto. La testa si inarcò sul cuscino. Stava avendo qualcosa
di simile a un attacco epilettico, per quanto breve: finì prima che
riuscissi ad avvicinarmi al letto. Rantolò e fece un respiro profondo
per riempire i polmoni d'aria. «Scusa, mi dispiace\ldots»

«Non ti scusare.»

«Non riesco a controllarlo, scusami.»

«Lo so che non puoi. Va tutto bene, Jase.»

«Non incolpare loro per ciò che mi sta capitando.»

«Loro chi - gli Ipotetici?»

Provò a sorridere, anche se stava chiaramente soffrendo.

«Dobbiamo trovargli un nome nuovo, ti va? Non sono poi così ipotetici
come una volta. Ma non incolparli. Non sanno ciò che mi succede. Mi
trovo sotto la loro soglia di astrazione.»

«Non so cosa vuol dire.»

Parlò veloce e con entusiasmo, come se parlare fosse una piacevole
distrazione dalla sofferenza fisica. Oppure un altro sintomo. «Tu e io,
Tyler, siamo colonie di cellule viventi, d'accordo? E se tu danneggiassi
un numero sufficiente delle mie cellule, io morirei e tu mi avresti
assassinato. Ma se ci stringiamo la mano e perdo qualche cellula cutanea
durante la procedura nessuno di noi si accorgerebbe minimamente della
perdita. Sarebbe invisibile. Viviamo a un certo livello di astrazione;
interagiamo come corpi, non come colonie di cellule. La stessa cosa vale
per gli Ipotetici. Popolano una zona di universo più vasta rispetto a
noi.»

«Questo li giustifica a uccidere le persone?»

«Sto parlando della loro percezione, non della loro moralità. La morte
di ogni singolo essere umano - la \emph{mia} morte - potrebbe essere
significativa per loro, se potessero inserirla nel giusto contesto. Ma
non possono.»

«Lo hanno fatto prima, però, hanno creato altri mondi-Spin - non è forse
una delle cose che i replicatori avevano scoperto prima che gli
Ipotetici li spegnessero?»

«Altri mondi-Spin. Sì. Molti. La rete degli Ipotetici si è sviluppata
fino ad abbracciare gran parte delle zone abitabili della galassia, e
questo è ciò che fanno quando s'imbattono in un pianeta che ospita una
specie senziente dotata di un certo grado di maturità e che riesce a
usare degli strumenti: lo rinchiudono in una membrana dello Spin.»

Immaginai dei ragni che avvolgono le loro vittime nella seta. «Perché,
Jase?»

La porta si aprì. Carol era tornata, portando un lumino su un piccolo
piatto di porcellana. Lo poggiò sulla credenza e lo accese con un
fiammifero. La fiamma danzò, messa a repentaglio dalla brezza che
arrivava dalla finestra.

«Per preservarla» disse Jason.

«Preservarla da cosa?»

«Dall'invecchiamento e dalla probabile morte. Le culture tecnologiche
sono mortali, come tutto il resto. Prosperano fino all'esaurimento delle
loro risorse; quindi muoiono.»

``Ammesso che succeda'', pensai. ``Ammesso che non continuino a
prosperare, a espandersi oltre i loro sistemi solari, trapiantandosi
sulle stelle\ldots''

Ma Jason aveva anticipato la mia obiezione. «I viaggi interstellari a
velocità subluce sono lenti e inefficaci per chi vive quanto un essere
umano. Forse avremmo potuto essere un'eccezione alla regola. Ma gli
Ipotetici esistono da moltissimo tempo. Prima di mettere a punto la
membrana dello Spin, hanno visto innumerevoli mondi abitati annegare nei
loro stessi effluvi.»

Fece un respiro e sembrò strozzarsi. Carol si girò verso di lui. La sua
maschera di professionalità scivolò via, e che Jason impiegò per
riprendersi riprendersi sembrò palesemente terrorizzata: non un medico,
ma una donna col figlio morente.

Jase, per fortuna, non vide nulla. Inghiottì a fatica e ricominciò a
respirare in modo regolare.

«Ma perché lo Spin, Jase? Ci spinge verso il futuro, ma non cambia
niente.»

«Al contrario» ribatté. «Cambia tutto.»

ooo

Il paradosso dell'ultima notte di Jason fu che il suo discorso diventò
sempre più strano e intermittente sebbene la conoscenza da lui acquisita
sembrasse espandersi in maniera esponenziale. Credo che in quelle poche
ore abbia imparato molto più di quanto potesse condividere, e ciò che
condivise fu memorabile - travolgente nel suo potere esplicativo e
stimolante per ciò che lasciava intendere sul destino dell'umanità.

Lasciando perdere il trauma e il suo tentativo angosciante di trovare le
parole giuste, ciò che disse fu\ldots{}

Be', iniziò con: «Prova a vedere le cose dal loro punto di vista. Gli
Ipotetici - sia che li consideriamo uno o molti organismi - si sono
evoluti da alcuni primitivi dispositivi di von Neumann per occupare la
nostra galassia. L'origine di quei primi dispositivi è avvolta
nell'oscurità. I loro discendenti non ne hanno memoria diretta, allo
stesso modo in cui tu e io non abbiamo un ricordo dell'evoluzione umana.
Forse erano il prodotto di una cultura biologica che ci ha preceduto e
di cui non rimane traccia; o magari erano migrati da una galassia più
vecchia. In entrambi i casi, gli Ipotetici di oggi appartengono a una
stirpe più antica di quanto si possa immaginare.

«Hanno visto un numero sterminato di specie senzienti evolversi ed
estinguersi su pianeti come il nostro. Potrebbero anche aver gettato il
seme di diversi processi evolutivi, trasportando passivamente materiale
organico da una stella all'altra. E sono stati spettatori di culture
biologiche che hanno creato rudimentali reti di von Neumann come
sottoprodotto della loro crescente (ma in fondo insostenibile)
complessità - non una ma svariate volte. Per gli Ipotetici queste
culture erano come vivai di replicatori: strane, prolifiche e fragili.

«Dal loro punto di vista questa gestazione, incessante ma a singhiozzi,
di elementari reti di von Neumann, seguita dal rapido collasso dei
pianeti d'origine, era insieme un mistero e una tragedia. Un mistero
perché gli eventi transitori, a livello biologico, erano per loro
difficili da capire o addirittura da cogliere. Una tragedia perché
avevano iniziato a concepire queste culture progenitrici come reti
biologiche fallimentari, simili a loro stessi quando erano in crescita
verso una reale complessità, ma annientati prematuramente da ecosistemi
planetari limitati.

«Secondo gli Ipotetici, quindi, lo Spin avrebbe dovuto proteggere noi e
decine di civiltà affini, emerse su altri mondi sia prima che dopo,
nella parte migliore della nostra esistenza tecnologica. Ma senza
considerarci pezzi da museo, congelati sul posto per essere mostrati al
pubblico. Gli Ipotetici hanno ridefinito il nostro destino. Ci hanno
sospesi in un tempo rallentato mentre mettevano insieme i pezzi di un
sontuoso esperimento, un esperimento formulato nel corso di miliardi di
anni che sta volgendo all'obiettivo finale: la costruzione di un
paesaggio biologico molto espanso in cui le nostre culture, altrimenti
spacciate, possano espandersi e finalmente incontrarsi e mischiarsi.»

ooo

Non ne afferrai immediatamente il significato.

«Un ambiente biologico espanso? Più grande della stessa Terra?»

Ci stavamo avvicinando al buio totale. Le parole di Jason venivano
interrotte da convulsioni e suoni involontari, che per scelta ho omesso
di raccontare. Di tanto in tanto gli controllavo il battito cardiaco,
rapido ma sempre più debole.

«Gli Ipotetici,» riprese, «possono manipolare il tempo e lo spazio. Le
prove sono tutte attorno a noi. Ma la creazione di una membrana
temporale non è né l'inizio né la fine delle loro capacità. Possono
connettere, attraverso anelli spaziali, il nostro pianeta ad altri
simili\ldots{} nuovi pianeti, alcuni progettati artificialmente e
accuditi, verso cui si può viaggiare \emph{in un attimo} e con
\emph{facilità} per mezzo di collegamenti, strutture e ponti assemblati
dagli Ipotetici partendo da - se è davvero possibile - materiale fatto
di stelle morte, stelle di neutroni\ldots{} strutture trascinate per lo
spazio, senza fretta, con pazienza, nel corso di milioni di anni\ldots»

Carol sedeva accanto a lui su un lato del letto e io sull'altro. Gli
tenevo ferme le spalle quando aveva le convulsioni e Carol gli
accarezzava la testa durante quegli intervalli in cui non riusciva a
parlare. I suoi occhi scintillavano al lume della candela e lui fissava
intensamente il nulla assoluto.

«La membrana dello Spin è ancora al suo posto, lavora, pensa, ma la
funzione temporale è finita, completata\ldots{} \emph{questo} erano gli
scintillii, il sottoprodotto di un processo di desintonizzazione, e ora
la membrana è stata resa permeabile in modo da farci passare attraverso
qualcosa, qualcosa di \emph{enorme}\ldots»

In seguito divenne chiaro ciò che voleva dire. Al tempo ero sbigottito e
sospettavo che avesse iniziato a soffrire di demenza, una sorta di
sovraccarico metaforico governato dalla parola ``rete''.

Mi sbagliavo, ovviamente.

\emph{Ars moriendi ars vivendi est}: l'arte di morire è l'arte di
vivere. Lo avevo letto da qualche parte, qualche giorno dopo la laurea e
mi tornò in mente mentre ero seduto al suo fianco. Jason morì come aveva
vissuto, all'eroica ricerca della conoscenza. Il suo dono al mondo
sarebbero stati i frutti di tale conoscenza, non accumulata da uno solo
ma distribuita liberamente.

Ma l'altro ricordo che mi saltò in mente, mentre la sostanza del sistema
nervoso di Jason veniva trasformata ed erosa dagli Ipotetici in una
maniera che forse non sapevano essere letale, fu di un pomeriggio di
tanto tempo prima, quando era montato in sella alla mia bici comprata al
mercatino dell'usato sfrecciando sulla Bantam Hill Road. Pensai a come
aveva controllato quel veicolo scassato in modo abile, facendo quasi un
balletto, fino a che non ne era rimasto più nulla tranne balistica e
velocità, nell'inevitabile collasso dell'ordine nel caos.

Il suo corpo da Quarto era una macchina perfettamente bilanciata. Non
morì con tranquillità. Poco prima di mezzanotte, Jason perse l'uso della
parola e fu in quel momento che cominciò a sembrare spaventato e non più
del tutto umano. Carol gli tenne la mano e gli sussurrò che era al
sicuro, che era a casa. Non so se quella consolazione lo raggiunse tra
le strane e intricate camere in cui era entrata la sua mente. Spero di
sì.

Poco dopo i suoi occhi rotearono all'indietro e i muscoli si
rilassarono. Il corpo lottò ancora, prendendo fiato in modo convulso
fino quasi al mattino.

Poi lo lasciai con Carol, che gli accarezzava la testa con infinita
dolcezza e gli parlava sottovoce come se potesse ancora sentirla, e non
mi accorsi che il Sole, quando sorse, non era più rigonfio e rosso ma
radioso e perfetto com'era prima della fine dello Spin.
